% DANH MỤC KÝ HIỆU VÀ CHỮ VIẾT TẮT: đặt trước Chương 1; mọi ký hiệu hoặc chữ viết tắt mới xuất hiện trong báo cáo được thêm vào bảng này.
\clearpage
\section*{DANH MỤC KÝ HIỆU VÀ CHỮ VIẾT TẮT}
\addcontentsline{toc}{section}{Danh mục ký hiệu và chữ viết tắt}

\begin{center}
\small
\renewcommand{\arraystretch}{1.25}
\begin{tabular}{@{}p{2.6cm}p{13.4cm}@{}}
\toprule
Ký hiệu & Ý nghĩa \\
\midrule
\multicolumn{2}{@{}l}{\textit{Bộ dữ liệu}} \\
HU & Hồ Huron: mực nước hồ hằng năm \\
LY & Linh miêu: logarit số linh miêu bị bẫy hằng năm \\
DR & DAX dạng suất sinh lợi: suất sinh lợi log hằng ngày của chỉ số DAX \\
NI & Nile: lưu lượng hằng năm của sông Nile \\
AP & AirPassengers: logarit số hành khách hàng không hằng tháng \\
DP & DAX dạng log giá: logarit giá đóng cửa hằng ngày của chỉ số DAX \\
\midrule
\multicolumn{2}{@{}l}{\textit{Mô hình và phương pháp}} \\
AR & tự hồi quy (autoregressive) \\
MA & trung bình trượt (moving average) \\
ARMA & tự hồi quy trung bình trượt \\
ARIMA & ARMA cho chuỗi đã lấy sai phân (tích hợp) \\
SARIMA & ARIMA có thành phần mùa vụ \\
GARCH & mô hình phương sai có điều kiện thay đổi theo thời gian \\
ACF, PACF & hàm tự tương quan, hàm tự tương quan riêng phần \\
ADF & kiểm định nghiệm đơn vị Dickey--Fuller mở rộng \\
KPSS & kiểm định tính dừng Kwiatkowski--Phillips--Schmidt--Shin \\
DM & kiểm định Diebold--Mariano \\
AIC, BIC & tiêu chuẩn thông tin Akaike, tiêu chuẩn thông tin Bayes \\
iid & độc lập và cùng phân phối \\
\midrule
\multicolumn{2}{@{}l}{\textit{Thước đo sai số}} \\
MAE, RMSE & sai số tuyệt đối trung bình, căn bậc hai của sai số bình phương trung bình \\
MAPE, MASE & sai số phần trăm tuyệt đối trung bình, sai số tuyệt đối trung bình đã chuẩn hóa \\
MCSE & sai số Monte Carlo của một kết quả mô phỏng \\
\midrule
\multicolumn{2}{@{}l}{\textit{Xếp loại mức độ phù hợp}} \\
P, M, K & phù hợp, phù hợp một phần, không phù hợp \\
\bottomrule
\end{tabular}
\end{center}
